% Generated draft: requires mathematical and source review.
\documentclass[11pt]{article}
\usepackage[margin=25mm]{geometry}
\usepackage{fontspec}
\setmainfont{Latin Modern Roman}
\usepackage{amsmath,amssymb,mathtools,unicode-math}
\setmathfont{Latin Modern Math}
\usepackage{enumitem,needspace,xurl,hyperref}
\setlength{\parindent}{0pt}
\setlength{\parskip}{5pt}
\setlength{\emergencystretch}{3em}
\newcommand{\N}{\mathbb{N}}
\newcommand{\Z}{\mathbb{Z}}
\newcommand{\Q}{\mathbb{Q}}
\newcommand{\R}{\mathbb{R}}
\newcommand{\C}{\mathbb{C}}
\newcommand{\F}{\mathbb{F}}
\newcommand{\A}{\mathbb{A}}
\newcommand{\PP}{\mathbb{P}}
\newcommand{\E}{\mathbb{E}}
\newcommand{\eps}{\varepsilon}
\newcommand{\dd}{\,\mathrm{d}}
\begin{document}
\section*{UnsolvedMath ID: 800910}
\textbf{Source problem code:} catalogue entry
\subsection{Definitions and mathematical statement}
Let \(\Omega\) denote the local complex domain intended in the source; its dimension, geometry, and boundary assumptions are not specified in the supplied excerpt. A plurisubharmonic function on \(\Omega\) is an upper semicontinuous function with values in \([-\infty,\infty)\) whose restriction to every complex line is subharmonic or identically \(-\infty\) on each connected component of the intersection. Write \(\operatorname{PSH}(\Omega)\) for this class, subject to the source's unspecified convention concerning functions identically \(-\infty\) on a connected component.

\textbf{Question 42 (attributed to Guedj).} In the local setting and with the complex Monge--Amp\`ere flow referred to in the source, can one run that flow from every initial function
\[
 u_0\in\operatorname{PSH}(\Omega),
\]
without requiring \(u_0\) to be bounded?

The supplied context states that the flows under discussion include a local version of the K\"ahler--Ricci flow as a particular case and can be run from bounded plurisubharmonic initial data. It motivates the question by the corresponding global results cited there as [GZ13, DiNL14b].

\textbf{Unresolved formulation.} The phrase ``the above complex Monge--Amp\`ere flow'' refers to material absent from the supplied record. No evolution equation, coefficients, boundary conditions, time interval, solution class, or meaning of attaining the initial datum is supplied. Consequently, the question above preserves the intended extension from bounded to arbitrary plurisubharmonic initial data, but does not specify a complete existence assertion. These missing conditions must be recovered from the source before a precise formulation can be finalized.
\subsection{Short English statement}
Guedj asks whether the local complex Monge--Amp\`ere flow discussed in the source, reportedly available for bounded plurisubharmonic initial data, can also start from an arbitrary plurisubharmonic function. The governing flow and the required solution and initial convergence notions are absent from the supplied excerpt, so the formulation remains incomplete.
\subsection{Sources}
\begin{itemize}
\item Supplied local UnsolvedMath registry record, permanent ID \(800910\), source problem code \texttt{catalogue entry}, titled ``Open Problems in Pluripotential Theory --- Question 42.'' The registry lists Dinew, Guedj, and Zeriahi as proposers, gives the year 2015, and labels the problem open. This status is supplied registry information, not an independently verified mathematical or literature assessment.
\item Dinew, Guedj, and Zeriahi, \emph{Open Problems in Pluripotential Theory} (2015), Question 42, attributed in the excerpt to Guedj. Supplied source URL: \url{https://arxiv.org/abs/1511.00705}. The record reports access on 2026-10-16 and extraction by \texttt{arxiv\_latex}; these are supplied provenance claims, not independent source verification.
\item The excerpt cites [GZ13] and [DiNL14b] for corresponding global results. Their bibliographic entries and precise results are not included in the supplied record.
\end{itemize}
\end{document}
